% Self-contained body; main.tex supplies amsmath, amssymb, amsthm and
% theorem/proposition/corollary/remark environments. No bespoke macros.
\section{An exact quadratic diagnostic and its limits}
\label{sec:quadratic-theory}

The results in this section are specializations of classical Gaussian
moment identities, spectral approximation, and stochastic approximation.
Their purpose is to distinguish population inversion, finite-sample
estimation, and actual empirical optimization. They do not establish a
substantially original learning algorithm or an advantage of autonomous
neural growth. All probability and expectation statements below are exact
under the stated model; the numerical checks in the experiments are
separate evidence.

Fix integers $d,n\geq1$. Matrices are real and symmetric unless indicated
otherwise, with inner product $\langle A,B\rangle=\operatorname{tr}(AB)$.
Write $\|\cdot\|_F$, $\|\cdot\|_{\mathrm{op}}$, and $\|\cdot\|_*$ for
Frobenius, operator, and nuclear norms. Set
\[
 X\sim N(0,I_d),\qquad H(X)=XX^\top-I_d,\qquad
 q_A(X)=\langle A,H(X)\rangle.
\]
The unknown fixed teacher is $A_*$. Labels obey
\begin{equation}
 Y=q_{A_*}(X)+\xi,\qquad
 \mathbb E\xi=0,\quad \mathbb E\xi^2=\sigma^2<\infty,
 \quad \xi\ \hbox{independent of }X.
 \label{eq:quadratic-model}
\end{equation}
Distinct examples have independent copies of $(X,\xi)$. Define population
squared risk $R(B)=\mathbb E[(Y-q_B(X))^2]$.

A signed centered quadratic network
\[
 f(x)=\sum_{j=1}^k a_j\{(w_j^\top x)^2-\|w_j\|_2^2\}
\]
represents $B=\sum_{j=1}^k a_jw_jw_j^\top$. Its minimum width is
$\operatorname{rank}(B)$: any such sum has rank at most $k$, while an
eigendecomposition represents $B$ with exactly its number of nonzero
eigenvalues. The zero predictor has width zero. A constant output can
implement the centering. Negative $a_j$ are allowed throughout; restricting
all outputs to be nonnegative changes the operation class to positive
semidefinite matrices and does not preserve every result below.

\subsection{Population geometry}

\begin{theorem}[T1: Gaussian quadratic geometry]
\label{thm:quadratic-geometry}
For symmetric $A,B$,
\begin{equation}
 \mathbb E[q_A(X)q_B(X)]=2\operatorname{tr}(AB).
 \label{eq:quadratic-isometry}
\end{equation}
Consequently, under \eqref{eq:quadratic-model},
\begin{equation}
 R(B)-R(A_*)=2\|B-A_*\|_F^2,
 \qquad \frac12\mathbb E[YH(X)]=A_*.
 \label{eq:quadratic-risk-moment}
\end{equation}
\end{theorem}
\begin{proof}
For standard Gaussian coordinates, second and fourth moments are
\[
 \mathbb E[X_iX_j]=\delta_{ij},\qquad
 \mathbb E[X_iX_jX_kX_l]
 =\delta_{ij}\delta_{kl}+\delta_{ik}\delta_{jl}
  +\delta_{il}\delta_{jk}.
\]
Subtracting the product of second moments yields
$\mathbb E[H_{ij}H_{kl}]=\delta_{ik}\delta_{jl}
+\delta_{il}\delta_{jk}$. Multiply by $A_{ij}B_{kl}$ and sum all indices.
The symmetry of $A,B$ gives \eqref{eq:quadratic-isometry}. Taking one
matrix to be an arbitrary symmetric test matrix gives
$\mathbb E[q_A(X)H(X)]=2A$.

Let $E=A_*-B$. Independence and zero mean of $\xi$ give
$\mathbb E[\xi q_E(X)]=0$ and $\mathbb E[\xi H(X)]=0$.
Thus $R(B)=\mathbb E[q_E(X)^2]+\sigma^2
=2\|E\|_F^2+\sigma^2$, while $R(A_*)=\sigma^2$.
The second identity follows by taking the moment of $Y$.
\end{proof}

Without centering, the corresponding quadratic prediction difference has
second moment $2\|A-B\|_F^2+\operatorname{tr}(A-B)^2$.
Thus centering is part of the model, not an inconsequential convention.

\subsection{What repeated growth from a fixed summary computes}

\begin{theorem}[T2: fixed-summary spectral growth]
\label{thm:spectral-collapse}
Fix a symmetric matrix $S$ and set $B_0=0$. At each nonterminal step,
select a unit eigenvector $v_t$ of $S-B_t$ whose eigenvalue $\lambda_t$
has maximum absolute value, and set
\[
 B_{t+1}=B_t+\lambda_tv_tv_t^\top.
\]
For a fixed $\tau\geq0$, terminate before this update if
$|\lambda_t|\leq\tau$. Every nonzero update selects a previously
unselected eigencomponent of $S$, in nonincreasing absolute-eigenvalue
order. In at most $d$ updates the terminal result is
\begin{equation}
 B_{\rm stop}=\sum_{j:\,|s_j|>\tau}s_ju_ju_j^\top,
 \qquad S=\sum_{j=1}^d s_ju_ju_j^\top.
 \label{eq:hard-truncation}
\end{equation}
Intermediate truncations can be nonunique at ties.
The final strict-threshold result is invariant to these choices.
\end{theorem}
\begin{proof}
We maintain an orthonormal set of selected eigenvectors of $S$ such that
$B_t$ is the sum of their original eigencomponents. This holds at $t=0$.
The residual $S-B_t$ is zero on their span and equals $S$ on its
orthogonal complement. If a nonzero residual eigenvalue remains, every
unit eigenvector with that nonzero eigenvalue is orthogonal to the
selected span. It is an eigenvector of $S$ with the same eigenvalue.
The update therefore extends the orthonormal selected set and removes
exactly that one component from the residual. Selecting maximum absolute
value establishes the order. At most $d$ nonzero components can be
removed. The stopping rule removes all and only components with
absolute eigenvalue strictly greater than $\tau$, proving
\eqref{eq:hard-truncation}.

Within an eigenspace with the same signed eigenvalue, any orthonormal
basis is permissible. If $s$ and $-s$ are tied in absolute value, one
may choose either eigenspace, but arbitrary mixtures across them need
not be eigenvectors and are not allowed by the algorithm. When the
residual is zero, the rule terminates, including when $\tau=0$.
\end{proof}

The directions in this construction can be learned from labels through
$S$. Nevertheless, after $S$ has been formed the trajectory is precisely
an eigendecomposition procedure. A one-shot estimator with the same
$S$, threshold and eigensolver produces the same predictor. This statement
does not identify the trajectory of a method that subsequently refits
neurons using empirical loss.

For an observed dataset, define
\begin{equation}
 S_n=\frac1{2n}\sum_{i=1}^nY_iH_i,
 \qquad
 T_n(B)=\frac1{2n}\sum_{i=1}^n\langle B,H_i\rangle H_i,
 \quad H_i=H(X_i).
 \label{eq:empirical-moment-operator}
\end{equation}
The actual residual moment satisfies the algebraic identity
\begin{equation}
 G_n(B):=\frac1{2n}\sum_{i=1}^n
       \{Y_i-q_B(X_i)\}H_i=S_n-T_n(B).
 \label{eq:actual-residual}
\end{equation}
It is the negative Euclidean matrix gradient of
$L_n(B)=(4n)^{-1}\sum_i\{Y_i-q_B(X_i)\}^2$.
The surrogate $S_n-B$ replaces the empirical operator $T_n$ by its
population mean. For fixed $B$ independent of the sample,
$\mathbb E[T_n(B)\mid B]=B$; equality on the realized sample does not
follow. If $B$ uses the same sample, even this conditional identity can
fail.

\begin{remark}[An explicit adaptive-reuse counterexample]
\label{rem:adaptive-counterexample}
Take $d=n=1$, $A_*=1$, $\sigma=0$, and $H=X^2-1$. Then $Y=H$.
Choose $B=S_1=H^2/2$, which is a legitimate data-dependent estimate.
Equation~\eqref{eq:empirical-moment-operator} gives $T_1(B)=B^2$.
Thus $\mathbb E[T_1(B)\mid B]=B^2$, generally different from $B$.
A same-sample residual correction is $C=B+G_1(B)=2B-B^2$, so
$|C-1|^2=(B-1)^4$. The fresh-block conditional variance in the next
theorem would instead be $14(B-1)^2$ and does not apply here.
\end{remark}

For exact empirical residual evaluation at arbitrary $B$, one must apply
$T_n$, using replay, its stored operator, or an equivalent representation.
In orthonormal symmetric coordinates, $T_n$ is one half the empirical
Gram operator of the $d(d+1)/2$ centered quadratic features. Its entries
include fourth-order input moments. This is the ordinary regression
normal equation; it is not a lower bound on physical memory, and no
minimum-bit claim follows from it.

\subsection{Exact sampling variance and fresh corrections}

\begin{theorem}[T3: variance of a fresh residual correction]
\label{thm:residual-variance}
Define
\[
 c_d=\frac{d^2+9d+14}{2}.
\]
For any fixed symmetric $E$ and $Z_E=\tfrac12q_E(X)H(X)$,
\begin{equation}
 \mathbb E Z_E=E,\qquad
 \mathbb E\|Z_E-E\|_F^2
 =c_d\|E\|_F^2+2\operatorname{tr}(E)^2.
 \label{eq:exact-design-variance}
\end{equation}
Let $\mathcal F$ contain the complete history before a new block, and
let $B$ be a square-integrable symmetric $\mathcal F$-measurable matrix.
Conditional
on $\mathcal F$, the new $n$ observations are independent copies of
\eqref{eq:quadratic-model} and independent of $\mathcal F$. Set
\[
 E=A_*-B,\qquad
 C=B+\frac1{2n}\sum_{i=1}^n
             \{Y_i-q_B(X_i)\}H_i.
\]
Then
\begin{align}
 \mathbb E[C\mid\mathcal F]&=A_*,\\
 \mathbb E[\|C-A_*\|_F^2\mid\mathcal F]
 &=\frac{c_d\|E\|_F^2+2\operatorname{tr}(E)^2
                +\sigma^2d(d+1)/4}{n}.
 \label{eq:fresh-conditional-variance}
\end{align}
\end{theorem}
\begin{proof}
Write $s=\|X\|_2^2$ and $U=X/\|X\|_2$, with any definition at
$X=0$, a null event. The radius $s\sim\chi_d^2$ is independent of the
uniform direction $U$ on the unit sphere. Put
$\theta=\operatorname{tr}(E)$ and $v=\|E\|_F^2$.
Rotational symmetry, or the Gaussian fourth-moment identity applied
before dividing out the radius, gives
\[
 \mathbb E[U^\top EU]=\frac{\theta}{d},\qquad
 \mathbb E[(U^\top EU)^2]=\frac{\theta^2+2v}{d(d+2)}.
\]
Indeed, the second formula follows from
$\mathbb E[(X^\top EX)^2]=\theta^2+2v$ and
$\mathbb E[s^2]=d(d+2)$ by radial independence.
Since $q_E=sU^\top EU-\theta$,
\begin{equation}
 \mathbb E[q_E^2\mid s]
 =\frac{(\theta^2+2v)s^2}{d(d+2)}
       -\frac{2\theta^2s}{d}+\theta^2.
 \label{eq:conditional-radial-moment}
\end{equation}
The chi-square moments are
\[
 \mathbb E[s^j]=\prod_{l=0}^{j-1}(d+2l),\qquad j=1,2,3,4.
\]
Substitute these into \eqref{eq:conditional-radial-moment}, after
multiplying respectively by $1,s,s^2$, to obtain
\begin{align*}
 \mathbb E[q_E^2]&=2v,\\
 \mathbb E[q_E^2s]&=2(d+4)v,\\
 \mathbb E[q_E^2s^2]&=8\theta^2+2(d+4)(d+6)v.
\end{align*}
Also $\|H\|_F^2=s^2-2s+d$. Hence
\begin{align*}
 \mathbb E[q_E^2\|H\|_F^2]
 &=8\theta^2+2\{(d+4)(d+6)-2(d+4)+d\}v\\
 &=8\theta^2+2(d^2+9d+16)v.
\end{align*}
By Theorem~\ref{thm:quadratic-geometry}, $\mathbb E Z_E=E$.
The identity
$\mathbb E\|Z_E-E\|_F^2=\mathbb E\|Z_E\|_F^2-\|E\|_F^2$
now proves \eqref{eq:exact-design-variance}.

For the noisy single-example correction, let
$W=Z_E+\tfrac12\xi H$. Conditional on the past, $E$ is fixed.
Independence and zero mean of $\xi$ eliminate all cross terms between
$Z_E-E$ and $\xi H$. Furthermore,
\[
 \mathbb E\|H\|_F^2=d(d+2)-2d+d=d(d+1).
\]
Consequently $\mathbb E[W\mid\mathcal F]=E$ and
\[
 \mathbb E[\|W-E\|_F^2\mid\mathcal F]
 =c_d\|E\|_F^2+2\theta^2+\frac{\sigma^2d(d+1)}4.
\]
The $n$ conditionally independent centered matrices $W_i-E$ have zero
expected pairwise Frobenius inner products. Their average therefore
has exactly $1/n$ times this variance. Since
$C-A_*=n^{-1}\sum_i(W_i-E)$, both claims follow.
\end{proof}

\begin{corollary}[Contraction with genuinely fresh blocks]
\label{cor:fresh-contraction}
Initialize a deterministic finite $B_0$ and repeat the unprojected update
in Theorem~\ref{thm:residual-variance}, with block size $n$ and fresh
independent blocks. Set
\[
 \kappa=\frac{c_d+2d}{n},\qquad
 \nu=\frac{\sigma^2d(d+1)}{4n}.
\]
For all integers $T\geq0$,
\begin{equation}
 \mathbb E\|B_T-A_*\|_F^2
 \leq\kappa^T\|B_0-A_*\|_F^2
       +\nu\sum_{j=0}^{T-1}\kappa^j.
 \label{eq:contraction-sum}
\end{equation}
If $n>c_d+2d$, the sum equals
$(1-\kappa^T)/(1-\kappa)$ and $\kappa<1$.
If instead every update is replaced by a best symmetric rank-at-most-$r$
Frobenius approximation, with known $r\geq\operatorname{rank}(A_*)$,
the same statement holds with $\kappa$ replaced by $4\kappa$ and $\nu$
by $4\nu$. That sufficient contraction condition is
$n>4(c_d+2d)$.
\end{corollary}
\begin{proof}
Cauchy--Schwarz on the eigenvalues gives
$\operatorname{tr}(E)^2\leq d\|E\|_F^2$.
Apply this inequality in \eqref{eq:fresh-conditional-variance}, then
take expectations. The resulting recurrence is
$e_{t+1}\leq\kappa e_t+\nu$, where
$e_t=\mathbb E\|B_t-A_*\|_F^2$. The right side is finite by induction;
iterating the recurrence proves \eqref{eq:contraction-sum}.
For rank truncation, the feasibility of $A_*$ and best-approximation
property give
\[
 \|P_r(C)-A_*\|_F
 \leq\|P_r(C)-C\|_F+\|C-A_*\|_F
 \leq2\|C-A_*\|_F.
\]
Square and condition on the past to obtain the stated factors of four.
A fixed measurable tie convention suffices to define the projection.
\end{proof}

This result has an $O(d^2)$ sufficient block size and an unconstrained
output that can require $d$ signed neurons. A supplied rank is privileged
information. These facts do not establish an optimal $O(dr)$ sample
bound, selected width, or a computational advantage over fixed-width
training. The matrix update is stochastic gradient descent for the
matrix-coordinate loss in \eqref{eq:actual-residual}; it is not the same
update as gradient descent in neuron factors. The exact variance formula
also quantifies the familiar control-variate effect of subtracting a
previously fitted predictor before estimating a residual.

\subsection{A wrong surrogate can propose harmful growth at an exact fit}

\begin{theorem}[T4: false growth at the optimum]
\label{thm:false-growth}
Assume $\sigma=0$ and $A_*\ne0$. From any $n\geq1$ fresh observations
form $S_n$ in \eqref{eq:empirical-moment-operator}, and consider the exact
fit $B=A_*$. Put $D=S_n-A_*$ and let $(\lambda,v)$ be a top-absolute-value
unit eigenpair of $D$. Consider the signed neuron update
$B'=A_*+\lambda vv^\top$. Its true risk increase and its improvement in
the surrogate $\widetilde R_{S_n}(B)=2\|B-S_n\|_F^2$ coincide:
\begin{equation}
 R(B')-R(A_*)
 =\widetilde R_{S_n}(A_*)-\widetilde R_{S_n}(B')
 =2\lambda^2.
 \label{eq:false-growth-gain}
\end{equation}
This quantity is strictly positive almost surely. With
$V=c_d\|A_*\|_F^2+2\operatorname{tr}(A_*)^2$,
\begin{equation}
 \frac{2V}{dn}
 \leq\mathbb E[R(B')-R(A_*)]
 \leq\frac{2V}{n}.
 \label{eq:false-growth-expectation}
\end{equation}
In contrast, the actual empirical residual moment at $A_*$ is exactly
zero on every sample.
\end{theorem}
\begin{proof}
Since $Y_i=q_{A_*}(X_i)$, every actual residual at $A_*$ is zero, proving
the last assertion. The population risk identity and
$\|vv^\top\|_F=1$ yield $R(B')-R(A_*)=2\lambda^2$.
Also $\langle D,vv^\top\rangle=v^\top Dv=\lambda$, so
\[
 \|D-\lambda vv^\top\|_F^2
 =\|D\|_F^2-2\lambda^2+\lambda^2
 =\|D\|_F^2-\lambda^2.
\]
This proves the surrogate identity.

Apply Theorem~\ref{thm:residual-variance} with $B=0$ and $\sigma=0$
to obtain $\mathbb E\|D\|_F^2=V/n>0$. At least one entry of $D$
therefore has positive second moment. Every entry is a polynomial in
the $nd$ independent Gaussian coordinates. At least one is a nonzero
polynomial; its zero set has Lebesgue measure zero. For completeness,
this property follows by induction on the number of variables: view a
nonzero multivariate polynomial as a polynomial in its last variable;
outside the common zero set of its nonzero coefficient polynomials,
it has finitely many roots, and Fubini's theorem finishes the induction.
The joint Gaussian law has a density, so $\mathbb P(D=0)=0$.
Thus $|\lambda|=\|D\|_{\mathrm{op}}>0$ almost surely.
Finally,
$\|D\|_F^2/d\leq\|D\|_{\mathrm{op}}^2\leq\|D\|_F^2$;
expectations give \eqref{eq:false-growth-expectation}.
\end{proof}

The exact fit is an oracle diagnostic checkpoint, not a checkpoint claimed
to be learned by the experimental controller. The theorem concerns the
unrestricted positive-score surrogate rule, with signed additions and
without a nonzero uncertainty threshold. It does not assert that a
thresholded or independently audited rule must accept. Any such accepted
nonzero update at the exact fit is harmful, but its acceptance probability
and expected accepted loss require a separate analysis. The theorem is
not a lower bound for every algorithm using $S_n$.

\subsection{A conditional width and risk certificate}

\begin{theorem}[Hard-threshold oracle inequality]
\label{thm:threshold-oracle}
Let $A_*,S$ be arbitrary symmetric matrices and let $b\geq0$ satisfy
$\|S-A_*\|_{\mathrm{op}}\leq b$. Let $s_1(A_*)\geq\cdots\geq s_d(A_*)$
denote the singular values of $A_*$, and let $\widehat A$ retain exactly
the eigencomponents of $S$ with absolute eigenvalue strictly greater
than $2b$. Then
\begin{align}
 \operatorname{rank}(\widehat A)
 &\leq\#\{j:s_j(A_*)>b\},
 \label{eq:threshold-rank}\\
 \|\widehat A-A_*\|_F^2
 &\leq\min_{0\leq r\leq d}
 \left\{3\sum_{j>r}s_j(A_*)^2+16b^2r\right\}.
 \label{eq:threshold-risk}
\end{align}
The corresponding excess prediction risk under
\eqref{eq:quadratic-model} is twice the bound in
\eqref{eq:threshold-risk}.
\end{theorem}
\begin{proof}
The singular-value perturbation bound
$|s_j(S)-s_j(A_*)|\leq\|S-A_*\|_{\mathrm{op}}$ implies that
$s_j(S)>2b$ entails $s_j(A_*)>b$, proving
\eqref{eq:threshold-rank}. One way to see this perturbation bound is to
apply the variational characterization
$s_j(M)=\min_{\dim W=d-j+1}\max_{u\in W,\|u\|=1}\|Mu\|$
and the triangle inequality, then interchange $S$ and $A_*$.

Hard thresholding at $2b$ minimizes, over symmetric $B$,
\[
 \|S-B\|_F^2+4b^2\operatorname{rank}(B).
\]
To justify this without a separate approximation theorem, let $B$ have
rank at most $r$ and let $P$ be the orthogonal projector onto its column
space. Then
\[
 \|S-B\|_F^2\geq\|(I-P)S\|_F^2
 =\operatorname{tr}(S^2)-\operatorname{tr}(PS^2).
\]
In an eigenbasis of $S^2$ ordered by decreasing eigenvalue, write
$p_j=u_j^\top Pu_j$. Since $0\leq p_j\leq1$ and
$\sum_jp_j=\operatorname{rank}(P)\leq r$,
$\operatorname{tr}(PS^2)\leq\sum_{j=1}^r s_j(S)^2$.
Retaining the largest $r$ singular components attains this bound and is
realizable as a symmetric eigencomponent sum. For each component,
retaining its value costs $4b^2$, whereas omitting it costs its squared
eigenvalue. At equality both choices minimize the objective, and the
theorem specifies omission.

Fix $r\in\{0,\ldots,d\}$ and let $A_{*,r}$ be a best symmetric
rank-at-most-$r$ approximation to $A_*$. Put
$k=\operatorname{rank}(\widehat A)$,
$e=\|\widehat A-A_*\|_F$, and
$a=\|A_{*,r}-A_*\|_F$. Comparing penalized objectives and expanding
around $A_*$ gives
\[
 e^2\leq a^2+2\langle S-A_*,\widehat A-A_{*,r}\rangle
                  +4b^2(r-k).
\]
The use of $r$ rather than the possibly smaller rank of $A_{*,r}$ can
only enlarge the right side. Operator/nuclear duality, the inequality
$\|M\|_*\leq\sqrt{\operatorname{rank}(M)}\|M\|_F$, and the triangle
inequality imply
\[
 e^2\leq a^2+2b\sqrt{k+r}(e+a)+4b^2(r-k).
\]
Apply $2b\sqrt{k+r}e\leq e^2/2+2b^2(k+r)$ and
$2b\sqrt{k+r}a\leq a^2/2+2b^2(k+r)$. Rearranging yields
$e^2\leq3a^2+16b^2r$. Finally,
$a^2=\sum_{j>r}s_j(A_*)^2$, and $r$ was arbitrary.
The prediction statement follows from
Theorem~\ref{thm:quadratic-geometry}. These arguments also cover $b=0$;
in that case $S=A_*$ and $\widehat A=A_*$ directly.
\end{proof}

The strict inequality in the threshold is necessary for the stated
width bound: for $d=1$, $b=1$, $A_*=1$, and $S=2$, retaining the component
at equality gives rank one, while the right side of
\eqref{eq:threshold-rank} is zero. The risk bound alone allows either
choice at equality because both minimize the penalized objective.

An explicit, conservative finite-sample radius follows without assuming
an optimal concentration theorem. For the raw moment $S_n$, let
\[
 V_* = c_d\|A_*\|_F^2+2\operatorname{tr}(A_*)^2
                +\frac{\sigma^2d(d+1)}4.
\]
Theorem~\ref{thm:residual-variance} gives
$\mathbb E\|S_n-A_*\|_F^2=V_*/n$. For any deterministic known
$\overline V\geq V_*$ and $\delta\in(0,1)$, take
$b=\sqrt{\overline V/(n\delta)}$. Markov's inequality and
$\|M\|_{\mathrm{op}}\leq\|M\|_F$ show that the assumptions of
Theorem~\ref{thm:threshold-oracle} hold with probability at least
$1-\delta$. If $\overline V=0$, the moment is exact almost surely.
Unknown teacher norms cannot simply be substituted by same-sample point
estimates without additional uncertainty control. This radius is
conservative and signal dependent; it does not establish the optimal
noisy low-rank rate. The oracle inequality is standard hard-threshold
reasoning and is included to make that limitation explicit.
